\section{Introduction}
\label{sec:introduction}

Stunting, defined as a height-for-age $z$-score more than two standard
deviations below the median of the WHO Child Growth Standards, remains one of
the most widely used indicators of chronic undernutrition in children under five
years of age \citep{who2006growth}. National programmes increasingly allocate
resources at the level of administrative districts, which creates a demand for
district-level prevalence estimates that are both timely and precise. The
household surveys that provide the underlying anthropometric data, however, are
designed to deliver reliable estimates at the national or provincial level. When
the same data are disaggregated to districts, many districts contain only a few
dozen measured children and the resulting direct estimates are too variable to
guide policy \citep{rao2015small}.

Small area estimation addresses this problem by combining the direct survey
estimates with auxiliary information that is available for all districts, such
as census aggregates, administrative registers or remotely sensed indicators.
Among the many available approaches, the area-level model of
\citet{fay1979estimates} is the most widely used in official statistics because
it requires only district-level direct estimates and their design-based
variances, and therefore respects the complex sampling design without requiring
access to unit-level records \citep{pfeffermann2013new}. The Fay--Herriot model
links the direct estimates to a linear regression on district covariates and
produces an empirical best linear unbiased predictor (EBLUP) that is a weighted
average of the direct estimate and the regression-synthetic estimate.

The linear mean is both the strength and the weakness of the Fay--Herriot model.
When the relationship between the outcome and a covariate is nonlinear, the
synthetic component is biased, and because the EBLUP shrinks most strongly in
districts with small samples, this bias is transferred precisely to the
districts for which the model is supposed to help most. In nutrition
applications the relationship between stunting prevalence and measures of
remoteness or market access is typically nonlinear: prevalence rises steeply
over the first few hours of travel time to the nearest town and then levels off
\citep{alderman2012access,headey2018drivers}. Transformations of the covariate
can partly absorb such patterns, but the choice of transformation is ad hoc and
rarely reported.

Several authors have relaxed the linearity assumption in small area models.
\citet{opsomer2008nonparametric} embedded penalized splines in a unit-level
nested error model and represented the spline coefficients as additional random
effects, which allows estimation with standard mixed model software.
\citet{ugarte2009spline} used a similar representation for disease mapping, and
\citet{sperlich2016model} discussed the trade-off between flexibility and
stability in model-based small area estimators. For area-level models, where the
number of districts rather than the number of individuals determines the
information available for the mean function, the use of splines has received
less attention; \citet{giusti2012semiparametric} and
\citet{marchetti2023robust} considered related ideas, but neither provides a
mean squared error estimator that accounts for the estimation of the smoothing
parameter.

In this paper we propose a spline-augmented Fay--Herriot (SAFH) model in which a
single continuous covariate, here a remoteness index, enters through a penalized
spline, while the remaining covariates enter linearly. Following
\citet{ruppert2003semiparametric}, the spline is written as a mixed model with a
second variance component, so that the smoothing parameter and the area-level
variance can be estimated jointly by restricted maximum likelihood (REML). We
derive the best predictor under the model, describe a parametric bootstrap for
its mean squared error (MSE) in the spirit of \citet{butar2003measuring}, and
compare the approach with the standard Fay--Herriot EBLUP and with a
Fay--Herriot model using a log-transformed covariate.

The methodology is motivated by the Northern Provinces Household Nutrition
Survey (NPHNS), a stratified two-stage cluster survey conducted in 2019 to
monitor the nutritional status of children in four northern provinces. The
survey was designed to produce provincial estimates, but district-level
estimates have been requested by the provincial health authorities to target a
new supplementary feeding programme. Our contribution is therefore both
methodological and applied: we show that modest nonlinearity in a single
covariate is enough to produce systematic bias in the standard estimator in
remote districts, and that the SAFH estimator removes most of this bias at a
small cost in variance.

The remainder of the paper is organized as follows. Section~\ref{sec:methods}
introduces the model, the estimation procedure and the MSE estimator.
Section~\ref{sec:simulation} reports a simulation study, and
Section~\ref{sec:application} presents the application to the NPHNS data.
Section~\ref{sec:discussion} concludes with a discussion of limitations and
extensions, and technical details of the bootstrap are given in
Appendix~\ref{app:bootstrap}.
